% =============================================================================
% Morehouse College Research Poster Template
% =============================================================================
% Brand colors:
%   Maroon  #840028 (Pantone 202 C)  ..  used as accent only
%   Gold    #C1A231 (Pantone 7753)   ..  used as accent only
%   Gray    #A7A8AA (Cool Gray 6)    ..  rules and secondary text
%   White   #FFFFFF                  ..  background (saves print ink)
%
% Compile with: pdflatex poster_morehouse.tex
% (run twice if any internal references need to resolve)
%
% Default size: A0 landscape (1189mm x 841mm). To change to US 48x36 inch
% landscape, replace "a0paper" below with "custom" and add the custom size
% lines shown in the comment block.
%
% To swap content: search for "EXAMPLE" comments throughout. Each \block
% command takes a title and a body. Add or remove blocks as needed.
% =============================================================================

\documentclass[25pt, a0paper, landscape, margin=0mm, innermargin=15mm,
               blockverticalspace=15mm, colspace=15mm,
               subcolspace=8mm]{tikzposter}

% For US 48 x 36 inch landscape instead, comment the line above and use:
%   \documentclass[25pt, custom, landscape, margin=0mm, innermargin=15mm,
%                  blockverticalspace=15mm, colspace=15mm]{tikzposter}
%   \geometry{paperwidth=48in, paperheight=36in}

\usepackage[utf8]{inputenc}
\usepackage[T1]{fontenc}
\usepackage{lmodern}
\usepackage{microtype}
\usepackage{graphicx}
\usepackage{booktabs}
\usepackage{array}
\usepackage{xcolor}
\usepackage{hyperref}
\hypersetup{colorlinks=true, linkcolor=morehousemaroon,
            urlcolor=morehousemaroon, citecolor=morehousemaroon}

% =============================================================================
% Morehouse brand colors
% =============================================================================
\definecolor{morehousemaroon}{HTML}{840028}
\definecolor{morehousegold}{HTML}{C1A231}
\definecolor{morehousegray}{HTML}{A7A8AA}
\definecolor{morehouseblack}{HTML}{000000}

% =============================================================================
% Custom Morehouse poster theme: white background, maroon and gold accents
% =============================================================================
\definebackgroundstyle{morehouseBg}{
    \draw[fill=white, draw=none] (bottomleft) rectangle (topright);
}
\definetitlestyle{morehouseTitle}{
    width=\paperwidth, roundedcorners=0, linewidth=0pt, innersep=12mm,
    titletotopverticalspace=0mm, titletoblockverticalspace=15mm
}{
    \begin{scope}[line width=0pt]
    \draw[draw=none, fill=white]
        (\titleposleft, \titleposbottom) rectangle (\titleposright, \titapostop);
    \draw[draw=none, fill=morehousemaroon]
        (\titleposleft, \titleposbottom)
        rectangle (\titleposright, \titleposbottom + 6mm);
    \draw[draw=none, fill=morehousegold]
        (\titleposleft, \titleposbottom + 6mm)
        rectangle (\titleposright, \titleposbottom + 8mm);
    \end{scope}
}
\definecolorpalette{morehousePalette}{
    \definecolor{colorOne}{named}{morehousemaroon}
    \definecolor{colorTwo}{named}{morehousegold}
    \definecolor{colorThree}{named}{morehousegray}
}
\definecolorstyle{morehouseColors}{\colorlet{notebgcolor}{morehousegold!30}}{
    \colorlet{backgroundcolor}{white}
    \colorlet{framecolor}{morehousemaroon}
    \colorlet{titlefgcolor}{morehousemaroon}
    \colorlet{titlebgcolor}{white}
    \colorlet{blocktitlebgcolor}{white}
    \colorlet{blocktitlefgcolor}{morehousemaroon}
    \colorlet{blockbodybgcolor}{white}
    \colorlet{blockbodyfgcolor}{morehouseblack}
    \colorlet{innerblocktitlebgcolor}{morehousegold}
    \colorlet{innerblocktitlefgcolor}{morehouseblack}
    \colorlet{innerblockbodybgcolor}{white}
    \colorlet{innerblockbodyfgcolor}{morehouseblack}
    \colorlet{notefgcolor}{morehouseblack}
    \colorlet{notebgcolor}{morehousegold!20}
    \colorlet{notefrcolor}{morehousegold}
}
\defineblockstyle{morehouseBlock}{
    titlewidthscale=1, bodywidthscale=1, titlecenter,
    titleoffsetx=0pt, titleoffsety=0pt, bodyoffsetx=0pt, bodyoffsety=0pt,
    bodyverticalshift=0mm, roundedcorners=0, linewidth=2pt,
    titleinnersep=8mm, bodyinnersep=8mm
}{
    \draw[color=morehousemaroon, fill=blockbodybgcolor,
          line width=\blocklinewidth] (blockbody.south west)
          rectangle (blockbody.north east);
    \ifBlockHasTitle
        \draw[draw=none, fill=morehousemaroon]
            (blocktitle.south west) rectangle (blocktitle.north east);
        \node[inner sep=0pt, anchor=center, text=white]
            at (blocktitle.center) {\textbf{\strut\blocktitletext}};
    \fi
}

\usebackgroundstyle{morehouseBg}
\usetitlestyle{morehouseTitle}
\usecolorpalette{morehousePalette}
\usecolorstyle{morehouseColors}
\useblockstyle{morehouseBlock}

% =============================================================================
% Title block content
% =============================================================================
\title{\parbox{0.95\linewidth}{\centering
    \color{morehousemaroon}\textbf{POSTER TITLE GOES HERE: Replace With Your Project's Working Title}
}}

\author{\parbox{0.95\linewidth}{\centering\color{morehouseblack}
    \large
    First Author$^{1}$, Second Author$^{1}$, Third Author$^{1}$, PI Name$^{1,2}$
}}

\institute{\parbox{0.95\linewidth}{\centering\color{morehouseblack}
    \normalsize
    $^{1}$Morehouse College, Atlanta, GA \quad
    $^{2}$Center for Broadening Participation in Computing
}}

% Optional logo. Drop your file into ../assets/ and uncomment.
% \titlegraphic{\includegraphics[height=4cm]{../assets/morehouse_logo.png}}

\begin{document}

\maketitle[width=\paperwidth, titletotopverticalspace=10mm]

% =============================================================================
% Row 1: Background and Dataset construction (two columns)
% =============================================================================
\begin{columns}

\column{0.5}
\block{Background}{
    \textbf{The problem.} Existing AI-generated image detection tools are
    trained on social media imagery: clean, well-lit, high-resolution.
    Real evidence footage in law enforcement settings is none of those
    things. Bodycam footage and surveillance video are typically grainy,
    compressed, low-resolution, and recorded under variable lighting.

    \vspace{4mm}
    \textbf{Why it matters.} Recent court cases have surfaced AI-generated
    images presented as evidence (Mendones v. Cushman, 2025). Federal
    Rule of Evidence 707 is being drafted to address admissibility of
    AI-generated content. Detection tools that fail on the very imagery
    courts must evaluate create a trust gap in the legal system.

    \vspace{4mm}
    \textbf{Our contribution.} We build the first domain-specific
    benchmark for evaluating AI-image detectors on law enforcement
    imagery, spanning multiple generator architectures and realistic
    image-quality conditions.
}

\column{0.5}
\block{Dataset Construction}{
    \begin{itemize}
        \item \textbf{Real images.} Stratified sample of N frames across
              14 categories of UCF Crime surveillance footage, plus N
              general-domain images from Open Images V7.
        \item \textbf{Synthetic images.} N images per generator, paired by
              prompt across three architectures:
              \begin{itemize}
                  \item FLUX.1-schnell (DiT, current SOTA)
                  \item Realistic Vision 5.1 (UNet, SD 1.5 fine-tune,
                        most-prevalent in-the-wild fabricator tool)
                  \item SDXL (UNet, established baseline)
              \end{itemize}
        \item \textbf{Quality levels.} Each image rendered at three
              degradation levels:
              \begin{itemize}
                  \item Clean: no degradation
                  \item Moderate: JPEG Q50, blur, contrast reduction
                  \item Heavy: JPEG Q30, downscaling, noise, blur
                        (simulates poor bodycam and old CCTV)
              \end{itemize}
        \item \textbf{Total:} N images at 3 quality levels = N evaluation
              instances.
    \end{itemize}
}

\end{columns}

% =============================================================================
% Row 2: Methods (full width, with subblocks)
% =============================================================================
\block{Methods}{
    \begin{minipage}[t]{0.48\linewidth}
        \textbf{Detector evaluation.} Each detection tool (commercial APIs
        and open-source models) is run against every image in the benchmark.
        We measure accuracy, precision, recall, F1, and AUC-ROC, broken
        down by:
        \begin{itemize}
            \item Real vs. synthetic
            \item Generator architecture (FLUX vs. RV vs. SDXL)
            \item Image category (surveillance, indoor, outdoor, etc.)
            \item Degradation level (clean, moderate, heavy)
        \end{itemize}
    \end{minipage}\hfill
    \begin{minipage}[t]{0.48\linewidth}
        \textbf{Paired-prompt design.} Each prompt is rendered through all
        three generators with matched seeds. This isolates detector
        sensitivity to architecture from sensitivity to content,
        enabling claims about generator-specific failure modes.

        \vspace{4mm}
        \textbf{Compute.} Generation runs on TACC Vista (NVIDIA GH200
        partition). Total GPU-hours: approximately N. Reproducibility
        package and full prompts are available in the project
        repository.
    \end{minipage}
}

% =============================================================================
% Row 3: Results (two columns: numbers + figure)
% =============================================================================
\begin{columns}

\column{0.5}
\block{Preliminary Results}{
    \textbf{Pilot evaluation} of an off-the-shelf detector
    (\texttt{umm-maybe/AI-image-detector}, 422 images) shows accuracy
    falling below random chance as image quality degrades.

    \vspace{6mm}
    \begin{center}
    \begin{tabular}{lrrr}
        \toprule
        \textbf{Quality} & \textbf{Acc.} & \textbf{F1} & \textbf{AUC} \\
        \midrule
        Clean    & 44.5\% & 37.4\% & 0.402 \\
        Moderate & 45.1\% & 17.1\% & 0.381 \\
        Heavy    & 34.5\% & 18.2\% & 0.265 \\
        \bottomrule
    \end{tabular}
    \end{center}

    \vspace{6mm}
    \textbf{Key finding.} Detection accuracy drops as image quality
    drops, exactly the direction it should not go for forensic use.
    Full evaluation across 3 generators and N detectors is in progress
    on TACC Vista.
}

\column{0.5}
\block{Discussion}{
    \textbf{Architecture generalization gap.} Detectors trained
    primarily on UNet-based generator outputs (SD 1.x, SDXL) are
    expected to underperform on transformer-based generators (FLUX).
    The benchmark is designed to quantify this gap.

    \vspace{4mm}
    \textbf{Domain gap.} Detectors trained on social media imagery
    are expected to underperform on surveillance domain content.
    Combining real surveillance frames (UCF Crime) with general-domain
    images (Open Images) lets us isolate this effect.

    \vspace{4mm}
    \textbf{Compounding under degradation.} The two gaps above are
    expected to compound non-linearly under realistic image-quality
    degradation, with the worst case (DiT generator at heavy
    degradation) showing the largest detection failure.

    \vspace{4mm}
    \textbf{Practical implication.} Existing AI-detection tools should
    not be used as standalone forensic instruments in law enforcement
    contexts without further calibration on domain-matched data.
}

\end{columns}

% =============================================================================
% Row 4: Acknowledgments and references (full width)
% =============================================================================
\block[bodyoffsety=0mm, titleoffsety=0mm]{Acknowledgments and References}{
    \begin{minipage}[t]{0.55\linewidth}
        \footnotesize
        \textbf{Funding.} This work is supported by the National
        Organization of Black Law Enforcement Executives (NOBLE)
        through a research grant to Morehouse College.

        \textbf{Compute.} Texas Advanced Computing Center (TACC) Vista
        and Stampede3 systems, NSF ACCESS allocation TRA25001.

        \textbf{Data.} UCF Crime Dataset (Sultani et al., 2018);
        Open Images Dataset V7 (Kuznetsova et al., 2020).
    \end{minipage}\hfill
    \begin{minipage}[t]{0.40\linewidth}
        \footnotesize
        \textbf{Project repository:} \url{https://github.com/...}\\
        \textbf{Dataset card:} \url{https://huggingface.co/datasets/...}\\
        \textbf{Contact:} ashley.scruse@morehouse.edu

        \vspace{4mm}
        \textit{Code and full prompts release upon paper acceptance.}
    \end{minipage}
}

\end{document}
